% 【本文件写什么】impl 实现层：dummy
%   - 定位：占位/链接桩实现，保证默认 feature 可编译。
%   - crate 路径：os/components/wateros-abi/abi-impl/impl-dummy/src/
%   - 如何实现对应 api-v0 trait：关键类型、状态机、数据结构
%   - 算法或硬件交互流程（可用伪代码/流程图）
%   - 本 impl 启用的 Cargo [features] 与 arch feature（impl-riscv64 等）
%   - 与其它 impl 变体的差异、切换 feature 时的影响
%   - 性能/内存假设与已知限制
% 【事实来源】
%   - 上述 crate 源码与 Cargo.toml
%   - os/feature-tree.txt 中指向本 impl 的 feature 链

\paragraph{职责}

\code{wateros-abi-impl-dummy}为Cargo工作区依赖解析与编译连通性占位，不包含真实平台ABI语义。对外仅提供\code{add(u64, u64) -> u64}模板函数，供crate内单测烟测使用。

\paragraph{与聚合层关系}

不通过\code{wateros-abi}聚合\code{lib.rs}导出；根\code{Cargo.toml}的\code{api-v0} feature链引用\code{impl-dummy/api-v0}仅为满足工作区成员解析。运行时主线选用\code{impl-linux-generic64}，切换至dummy不影响已启用号表的内核镜像。

\paragraph{限制}

无\code{SyscallNumberTable}实现，不参与trap或syscall分发。若仅启用\code{impl-dummy}而无\code{impl-linux-generic64}，\code{ActiveSyscallNumberTable}别名不存在，依赖活跃号表的crate无法编译。
